\section{Related Work}\label{sec:related}

\paragraph{Lattice Boltzmann Methods for Multiphase Flows.}
The lattice Boltzmann method (LBM) has emerged as a powerful tool for simulating multiphase flows due to its simplicity in handling complex boundary conditions and natural parallelization. Several multiphase models exist: the pseudopotential model \citep{shan1993lattice}, the color-gradient model \citep{gunstensen1991lattice}, the free-energy model \citep{swift1995lattice}, and the phase-field model \citep{lee2005lattice}. Among these, phase-field approaches based on the Allen-Cahn equation \citep{fakhari2017improved} have gained popularity for their ability to handle high density ratios with good stability.

\paragraph{High Density Ratio LBM.}
Achieving stable simulations at high density ratios (e.g., water/air at 828:1) remains challenging. Early index-function models were limited to density ratios below 10. \citet{inamuro2004lattice} achieved density ratios of 1000 using a free-energy approach with Poisson pressure correction. \citet{lee2003pressure} developed a pressure-evolution model reaching density ratios of 1000. More recently, \citet{fakhari2017improved} proposed an improved locality-preserving AC-LBM that supports density ratios up to 1000 with good stability. Our work builds on this formulation with additional stabilization for curved boundaries.

\paragraph{Droplet Impact on Curved Surfaces.}
The study of droplet dynamics on curved substrates has attracted significant attention. \citet{liu2015symmetry} experimentally demonstrated asymmetric bouncing on \emph{Echeveria} leaves, showing 40\% contact time reduction due to momentum asymmetry. Their lattice Boltzmann simulations (using the Lee-Liu free-energy model) confirmed the mechanism but were limited to a single Weber number. \citet{xu2022droplet} studied droplet impact on triangular ridges using pseudopotential LBM, finding that contact time reduction depends on the interplay between Weber number and ridge geometry. \citet{xiao2025numerical} investigated droplet impact and freezing on supercooled curved surfaces using thermal LBM with a pseudopotential model, providing orthogonal experimental design analysis of We, wall temperature, contact angle, and curvature ratio.

\paragraph{Wetting Boundary Conditions on Curved Surfaces.}
Accurate wetting boundary conditions are critical for simulating droplet-surface interactions. \citet{fakhari2017diffuse} proposed a geometric wetting approach for curved boundaries in phase-field LBM. \citet{connington2013lattice} developed a wetting model for curved surfaces using the free-energy approach. Recent work by \citet{sashko2024phase} extended phase-field LBM wetting to 3D curved geometries with interpolation-based mitigation of staircase approximations. \citet{compact2025stencil} proposed a compact-stencil wetting BC for 3D curved surfaces, eliminating nonphysical droplet sliding and detachment.

\paragraph{Curved Wetting Boundary Conditions in Phase-Field LBM (2022--2026).}
Several recent works have developed wetting boundary conditions (WBCs) for curved surfaces in the phase-field LBM framework. \citet{huang2022simplified} proposed a simplified geometric WBC for the conservative Allen-Cahn LBM using a hyperbolic tangent profile projection, validated at contact angles $\theta_{eq} = 45^\circ$--$135^\circ$ on cylindrical surfaces. \citet{zhang2022wetting} developed two geometric-formulation WBCs for modified phase-field LBM with large density ratios on flat substrates. \citet{huang2023simplified} extended the simplified scheme to a broader range of curved geometries. \citet{sashko2024phase} and \citet{compact2025stencil} addressed 3D curved geometries with interpolation-based and compact-stencil approaches, respectively. \citet{wang2024wetting} developed a WBC for three-dimensional curved geometries in the color-gradient LBM framework. \citet{ezzatneshan2021studying} applied the Allen-Cahn LBM to droplet impingement on hydrophilic and hydrophobic curved surfaces.

Despite this progress, all existing curved WBCs are purely geometric and parameter-free. They assume that the geometric formulation exactly enforces the prescribed contact angle regardless of curvature or hydrophobicity. Consequently, their validation is limited to modest curvature ratios and contact angles ($\theta_{eq} \leq 135^\circ$, with most cases at $\theta_{eq} \leq 120^\circ$). None of these studies identify or address the competition between the AC sharpening term and the geometric wetting correction.

\paragraph{Over-relaxation Techniques in LBM.}
Over-relaxation of boundary gradients is a known technique in LBM. \citet{ding2007wetting} used extrapolation-based wetting BCs with relaxation factors for free-energy LBM. Extrapolation- and interpolation-based wetting treatments with tunable relaxation were also explored for free-energy and phase-field models. However, these works focus on flat or mildly curved surfaces and do not provide a systematic calibration or physical motivation grounded in the AC sharpening dynamics. Our work differs in three ways: (1) we identify the AC sharpening as the specific source of resistance on strongly curved hydrophobic surfaces, (2) we propose $\alpha_{geo}$ motivated by the AC-wetting balance (Eq.~\ref{eq:alpha_derivation}), and (3) we show that the optimal $\alpha_{geo}$ depends on $\theta_{eq}$, $R^*$, and $D_0/\xi$---a systematic calibration across regimes not previously available.

\paragraph{Gap: The Untested Superhydrophobic Regime.}

A systematic survey of wetting BC validation across the phase-field LBM literature reveals a critical gap. \Cref{tab:lit_validation} summarizes the maximum validated contact angle and geometry complexity for each major contribution.

\begin{table}[htbp]
\centering
\caption{Validation coverage of wetting boundary conditions in phase-field LBM literature. No prior work validates the geometric WBC at superhydrophobic contact angles ($\theta_{eq} > 150^\circ$) on curved surfaces.}
\label{tab:lit_validation}
\begin{tabular}{c|c|c|c|c}
\hline
Reference & Max $\theta_{eq}$ & Curved? & Density ratio & Correction \\
\hline
\citet{fakhari2017diffuse} & $\sim 120^\circ$ & Cylinder & 1000:1 & Cubic surface energy \\
\citet{huang2022simplified} & $135^\circ$ & Sphere, spheroid & 1000:1 & Geometric (tanh-based) \\
\citet{zhang2022wetting} & $140^\circ$ & Flat only & 1000:1 & Geometric (virtual nodes) \\
\citet{sashko2024phase} & $\sim 120^\circ$ & 3D arbitrary & 1000:1 & Interpolation-based \\
\citet{compact2025stencil} & $\sim 120^\circ$ & 3D arbitrary & 1000:1 & Compact-stencil \\
\hline
\textbf{This work} & \textbf{$\mathbf{90^\circ}$--$\mathbf{162^\circ}$} & \textbf{Cyl. ridge, concave, convex} & \textbf{828:1} & \textbf{Geometric amplification ($\alpha_{geo}$)} \\
\hline
\end{tabular}
\end{table}

At $\theta_{eq} \leq 140^\circ$, the wall-normal gradient required by the geometric BC scales as $|\cot\theta_{eq}| \lesssim 1.19$, and the AC sharpening resistance is moderate. Most prior works validate geometric WBCs within this range and report good agreement with prescribed contact angles (Table~\ref{tab:lit_validation}). To our knowledge, the regime of strongly superhydrophobic curved surfaces ($\theta_{eq} > 164^\circ$, $|\cot\theta_{eq}| > 3.5$) has not been systematically tested with geometric WBCs; our work targets this regime and provides a systematic analysis of the AC-wetting under-correction there.

In summary, despite progress, no existing study provides: (a) identification of the AC sharpening term as the source of geometric BC inaccuracy on curved surfaces, (b) a physically motivated correction factor supported by scaling analysis, or (c) a comprehensive We $\times$ $R^*$ phase diagram for droplet impact on curved surfaces via Allen-Cahn LBM. Our work addresses all three.
